\documentclass[conference]{IEEEtran}
\IEEEoverridecommandlockouts

% *** PACKAGES ***
\usepackage{times}
\usepackage{graphicx}
\usepackage{amsmath}
\usepackage{amssymb}
\usepackage{url}
\usepackage{booktabs}
\usepackage{multirow}
\usepackage{array}
\usepackage{cite}
\usepackage{hyperref}
\usepackage{caption}
\usepackage{subcaption}
\usepackage{enumitem}
\usepackage{float}
\usepackage{tabularx}
\usepackage{makecell}
\usepackage{siunitx}

% Float spacing tuning
\setlength{\floatsep}{8pt plus 2pt minus 2pt}
\setlength{\textfloatsep}{10pt plus 2pt minus 4pt}
\setlength{\intextsep}{8pt plus 2pt minus 2pt}
\renewcommand{\topfraction}{0.95}
\renewcommand{\bottomfraction}{0.5}
\renewcommand{\textfraction}{0.05}
\renewcommand{\floatpagefraction}{0.9}

% Fixed-width column types
\newcolumntype{C}[1]{>{\centering\arraybackslash}p{#1}}
\newcolumntype{L}[1]{>{\raggedright\arraybackslash}p{#1}}
\newcolumntype{R}[1]{>{\raggedleft\arraybackslash}p{#1}}

% -----------------------------------------------------------
\begin{document}
% -----------------------------------------------------------

\title{NutriDermAI: A Multimodal Image--Text Fusion Framework
for ABCDE-Aware Skin Lesion Analysis and Interactive Dermatology VQA}

\author{
    \IEEEauthorblockN{Akhil Kulal}
    \IEEEauthorblockA{Department of Computer Engineering\\
    Veermata Jijabai Technological Institute\\
    Email: arkulal\_m24@ce.vjti.ac.in}
    \and
    \IEEEauthorblockN{Prof. Manasi Kulkarni}
    \IEEEauthorblockA{Department of Computer Engineering\\
    Veermata Jijabai Technological Institute\\
    Email: mukulkarni@ce.vjti.ac.in}
}

\maketitle

% -----------------------------------------------------------
\begin{abstract}
% -----------------------------------------------------------
Melanoma accounts for over 57,000 deaths annually, yet most AI-based
dermatology tools are image-only classifiers without structured ABCDE
(Asymmetry, Border, Colour, Diameter, Evolution) computation or
multimodal fusion. This paper presents \emph{NutriDermAI}, the first
fully integrated dermatology decision support system that computes
clinically calibrated ABCDE scores through principled image--text
fusion, not image-only analysis. Key novelties: (1)~multimodal ABCDE
score computation combining dermatoscopic features with clinical text,
(2)~transparent Clinical Knowledge JSON for auditable clinical
reasoning, (3)~end-to-end pipeline from segmentation to interactive
VQA. Four core stages are fully implemented and validated: skin lesion
segmentation (Dice 0.8506, Boundary F1 0.8815), ROI cropping (98.4\%
coverage, 12.3 ms), hierarchical classification (94.7\% binary, 89.2\%
7-class accuracy), and clinical NLP for evolution scoring (NER F1
0.87--0.89, ECE 0.0201). End-to-end integration achieves 517 ms latency
on NVIDIA A100, clinically acceptable for batch processing. This work
provides a validated foundation for ABCDE-aware dermatology AI suitable
for prospective clinical validation and resource-limited deployment.
\end{abstract}

\begin{IEEEkeywords}
Dermatoscopy, Skin Lesion, ABCDE Rule, Multimodal Learning,
Explainable AI, Medical VQA, Medical Report Processing, Image--Text
Fusion, Dermatology Chatbot
\end{IEEEkeywords}

% ===========================================================
\section{Introduction}
% ===========================================================

Skin cancer, including melanoma and non-melanoma lesions, continues to
present a major public health burden, with melanoma alone responsible
for over 57,000 deaths annually worldwide~\cite{sung2021_global_cancer_stats}.
Globally, more than 1 million skin cancer diagnoses are made each year,
with incidence rising 1.5--3\% annually in developed nations. Early detection
substantially improves five-year survival rates from under 30\% for late-stage
to over 98\% for localised melanoma~\cite{siegel2023_cancer_statistics}, making
timely and accurate diagnosis critical.

Dermatologists rely on dermatoscopic examination (magnified, illuminated lesion
inspection) guided by the ABCDE rule---Asymmetry, Border irregularity, Colour
variation, Diameter, and Evolution---to stratify malignant risk. However, three
clinical challenges limit this approach: (i)~\emph{Specialist shortage}: fewer
than 10,000 dermatologists serve over 8 billion people globally, with severe
disparities in resource-limited regions; (ii)~\emph{Subjective assessment}:
even among specialists, inter-observer agreement on ABCDE scores is 
moderate~\cite{kittler2011_abcde_variability}, introducing diagnostic 
inconsistency; and (iii)~\emph{Limited contextualisation}: dermatoscopic 
images alone lack integration with longitudinal history, family background, 
and symptom progression---critical signals in clinical practice. Classical 
ABCD/ABCDE-based expert systems have demonstrated that structured feature 
assessment can approximate dermatologists' diagnostic reasoning~\cite{chatterjee2021_abcd_expert_system,
senan2021_abcd_analysis,tas2022_fast_abcd_tds,
mevorach2025_ai_vs_derm_abcde,jutte2024_generative_ABCDE}, yet most do not 
integrate images and text, nor provide transparent, clinician-aligned explanations.

Concurrently, deep learning has achieved high performance on
dermatoscopic segmentation and classification tasks using U-Net
variants, hybrid CNN--Transformer designs, and large-scale datasets
such as HAM10000, ISIC 2016--2024, BCN20000, and
MiLK-10K~\cite{ronneberger2015_unet,arcunet2024_arcunet,
scalefusionnet2025,escunet2025,mrpunet2025,metaunet2025,biadatu2024,
mucmnet2024,isic2016_2024_datasets,tschsandl2018_ham10000,
combalia2019_bcn20000,milk10k2021_milk10k}. More recently, dermatology
foundation models and vision--language models (VLMs) such as PanDerm,
MM-Skin, and Derm1M have demonstrated how large-scale multimodal
supervision improves zero-shot recognition and text grounding in
dermatology~\cite{pan2025_panderm,zeng2025_mmskin,yan2025_derm1m,
aboulmira2025_attentionmultiskin,multimodal2025_skin_cancer_prediction}.

However, a critical gap remains: no existing system computes structured
ABCDE scores via principled multimodal fusion. Most AI systems segment
lesions or classify them in isolation, without grounding in the clinical
ABCDE framework. Those that do compute ABCDE often use ad-hoc image-only
approaches, ignoring Evolution (E), which requires clinical history.
This architectural gap motivates NutriDermAI.

Medical NLP and medical VQA have evolved rapidly, leveraging clinical
corpora like MIMIC-III/IV and architectures such as BioBERT, medical
multimodal LLMs, and RL-trained VQA
systems~\cite{lee2020_biobert,johnson2016_mimic3,hyland2023_mimic4note,
clinicaltextreview2020t,biomedblip2024_naseem,li2022_semantic_vqa,
dsap2024_dynamic_prompting,medvlmr12025,rlvalidation2023_vqa,
medthink2025,medicalmmllmreview2024}. Dermatology-specific VQA
benchmarks are emerging, yet they typically lack end-to-end ABCDE
computation and structured imaging--text
fusion~\cite{he2020_pathvqa,lau2018_vqarad,yim2024_dermavqa,
healmqa2024,faggioli2025_multimodal_sls_vqa,irel2025_clipseg,
dermatologychatbot2024,casestudy2024_wizard_of_oz,chatgptderm2025,
amie2025,convagents2025_effectiveness}.

\textbf{Problem Statement and Research Opportunity.} Current
dermatology AI systems suffer from two critical limitations: (i)~most
are unimodal (image-only), computing classification scores without
grounding in the clinical ABCDE framework; and (ii)~those that do
compute ABCDE metrics typically use ad-hoc image-only approaches that
cannot incorporate Evolution (E), a key criterion requiring longitudinal
clinical history. This architectural gap has prevented deployment of
ABCDE-aware, multimodal systems that could provide clinicians with both
high accuracy and transparent, evidence-grounded explanations.

\textbf{Solution and Vision.} NutriDermAI is an integrated multimodal
framework designed to close this gap. The system computes structured,
calibrated ABCDE scores via fusion of image and text features,
producing a Clinical Knowledge JSON that captures all reasoning in an
auditable, clinician-aligned format. By integrating this structured
representation with interactive VQA and knowledge-grounded conversational
AI, NutriDermAI supports multi-turn clinical dialogue while minimising
hallucination. The vision is:
\begin{itemize}[leftmargin=*]
    \item compute explicit, clinically calibrated ABCDE values aligned with 
          dermatology practice guidelines (not ad-hoc black-box scores);
    \item integrate images, patient history, and queries through principled 
          multimodal fusion (not image-only analysis);
    \item produce structured JSON explanations with evidence rankings suitable 
          for clinical audits and physician handoff; and
    \item enable interactive, context-aware question answering grounded in 
          ABCDE features and verified medical evidence.
\end{itemize}

\textbf{Key Contributions.} This paper presents three core contributions to
ABCDE-aware dermatology AI:

\begin{enumerate}[leftmargin=*]
    \item \textbf{Multimodal ABCDE computation.} Unlike prior image-only
          systems, we compute clinically calibrated ABCD scores via principled
          fusion of image features with clinical text (Stage 6). This ensures
          the Evolution (E) score is properly grounded in longitudinal history,
          addressing the architectural gap in existing systems.
    \item \textbf{Transparent clinical reasoning via structured JSON.} All
          downstream processing uses a Clinical Knowledge JSON as the canonical
          representation, enabling end-to-end audit trails, independent stage
          evaluation, and clinician-interpretable explanations aligned with
          ABCDE practice standards.
    \item \textbf{Fully integrated pipeline with validated core stages.} We
          demonstrate four production-ready stages (SLS, SLRC, HC, MedProc)
          with strong empirical results on standard benchmarks, plus integrated
          end-to-end processing achieving 517 ms latency. This provides a
          validated foundation for prospective clinical studies and deployment
          in resource-limited settings.
\end{enumerate}

% ===========================================================
\section{Background and Related Work}
% ===========================================================

\subsection{Problem Statement and Research Gaps}

Current dermatology AI systems exhibit distinct limitations that motivate
NutriDermAI's design:

\begin{enumerate}[leftmargin=*]
    \item \textbf{ABCDE-aware systems are rare and inconsistent.} While
          the ABCDE rule is the clinical gold standard for melanoma risk
          assessment, most AI systems either ignore it entirely or compute
          it ad-hoc without principled fusion of image and text signals.
    \item \textbf{Multimodal fusion is under-developed in dermatology.}
          Image-only classifiers miss critical temporal context; text-only
          systems miss geometric features. Rigorous fusion architectures that
          weight both modalities appropriately are lacking.
    \item \textbf{VQA systems lack ABCDE grounding.} Existing medical VQA
          benchmarks (VQA-RAD, PathVQA, DermaVQA) do not integrate
          structured ABCDE evidence, leading to potential hallucination.
    \item \textbf{End-to-end pipeline integration is missing.} While individual
          components (segmentation, classification, NLP) are well-studied,
          few works demonstrate a complete, clinical-grade pipeline with
          structured outputs, audit trails, and clinician-aligned explanations.
    \item \textbf{Fairness and bias remain unaddressed in deployments.}
          Most published systems lack evaluation across diverse skin tones,
          potentially perpetuating healthcare disparities.
\end{enumerate}

\subsection{ABCDE Rule and Classical Decision Support}

The ABCD/ABCDE rule offers a structured approach to melanoma risk
assessment by quantifying asymmetry, border irregularity, color
variation, diameter, and evolution over time. This rule was originally
formalised to address the global shortage of dermatologists, where
primary care physicians with limited dermoscopy training must triage
suspicious lesions~\cite{dreremoula2022_derm_shortage}. Classical expert
systems compute scores from dermatoscopic images and combine them into
total dermatoscopic scores (TDS)~\cite{chatterjee2021_abcd_expert_system,
senan2021_abcd_analysis,tas2022_fast_abcd_tds}. Recent work shows that
AI systems explicitly incorporating ABCDE analysis can approach
dermatologist-level decision-making~\cite{jutte2024_generative_ABCDE,
mevorach2025_ai_vs_derm_abcde}. Hybrid CNN+ABCD models further
demonstrate that explicit ABCD descriptors boost classification
performance when combined with deep features~\cite{shetty2022_hybridcnn_abcd}.

\subsection{Skin Lesion Segmentation}

Lesion segmentation is a prerequisite for ABCDE feature computation
and ROI cropping. U-Net and its variants remain dominant in medical
image segmentation~\cite{ronneberger2015_unet,reviewunet2024_unetreview}.
Dermatology-specialised architectures---ARC-UNet, ScaleFusionNet,
ESC-UNet, MRP-UNet, Meta-UNet, BiADATU-Net, MUCM-Net, and dual
CNN--ViT hybrids---have improved boundary precision and robustness via
residual connections, deformable attention, and multiscale
inputs~\cite{arcunet2024_arcunet,scalefusionnet2025,escunet2025,
mrpunet2025,metaunet2025,biadatu2024,mucmnet2024,hybriddualencoder2024}.
Multi-scale Swin Transformer models further improve segmentation across
diverse medical imaging
domains~\cite{swinmedical2025,attentionreview2023}.

\subsection{Lesion Classification and Foundation Models}

Large datasets---ISIC 2016--2024, HAM10000, BCN20000, and
MiLK-10K---have enabled high-performing classifiers and foundation
models~\cite{isic2016_2024_datasets,tschsandl2018_ham10000,
combalia2019_bcn20000,milk10k2021_milk10k}. EfficientFormerV2 provides
competitive accuracy with reduced memory
footprint~\cite{yu2022_efficientformerv2}, making it suitable for
multi-stage pipelines. PanDerm, MM-Skin, and Derm1M align dermatoscopic
images with rich textual descriptions and hierarchical
ontologies~\cite{pan2025_panderm,zeng2025_mmskin,yan2025_derm1m,
aboulmira2025_attentionmultiskin}, motivating shared multimodal encoders.

\subsection{Clinical NLP and Medical Text Understanding}

BioBERT and related biomedical encoders have improved named entity
recognition and classification over clinical
notes~\cite{lee2020_biobert}. MIMIC-III and MIMIC-IV provide rich
free-text and structured data for training and
evaluation~\cite{johnson2016_mimic3,hyland2023_mimic4note}. Systematic
reviews of clinical text data highlight both opportunities and pitfalls
for machine learning on heterogeneous EHR
text~\cite{clinicaltextreview2020t}. Retrieval-Augmented Generation
(RAG) over clinical knowledge bases has further been shown to reduce
hallucinations and ground LLM responses in verified medical
evidence~\cite{lewis2020_rag,xiong2024_benchmarking_rag_medical},
directly motivating the FAISS-based retriever in NutriDermAI's Stage~8.
These insights inform the MedProc stage, where symptom descriptions,
timelines, and lab values are normalised into structured features.

\subsection{Medical VQA and Multimodal LLMs}

Medical VQA benchmarks include VQA-RAD, PathVQA, TM-PathVQA, HEALMQA,
and dermatology-specific DermaVQA~\cite{lau2018_vqarad,he2020_pathvqa,
rajkhowa2024_tmpathvqa,yim2024_dermavqa,healmqa2024}. Architectures
range from BLIP-style encoders to semantic-attention VQA
layers~\cite{biomedblip2024_naseem,li2022_semantic_vqa}. Dynamic
prompt adaptation and RL-based optimisation reduce hallucinations and
improve reasoning~\cite{dsap2024_dynamic_prompting,medvlmr12025,
rlvalidation2023_vqa,medthink2025,madavan2025_medgrim}.
Dermatology-focused works demonstrate CLIPSeg-based multimodal fusion
and knowledge-enhanced
ensembles~\cite{faggioli2025_multimodal_sls_vqa,irel2025_clipseg,
h3n1lab2025_dermkem,lozano2025_synth_vqa,sviridov2025_multiagent_dialogue}.
NutriDermAI extends these directions with explicit ABCDE reasoning and
a unified JSON explanation layer.

\subsection{Dermatology Chatbots and Conversational AI}

Dermatology chatbots combining CNN-based lesion analysis with
rule-based triage logic have been
explored~\cite{dermatologychatbot2024,casestudy2024_wizard_of_oz}.
AMIE evaluates conversational diagnostic quality across specialties and
emphasises safety, uncertainty communication, and escalation
strategies~\cite{amie2025}. Reviews of LLMs in dermatology highlight
the need for guardrails and rigorous
evaluation~\cite{chatgptderm2025,medicalmmllmreview2024}. Studies on
conversational agents report improved patient engagement alongside
responsibility and ethical
concerns~\cite{convagents2025_effectiveness}. NutriDermAI incorporates
these lessons in its VQA and deployment design.

\subsection{Integration and System Design Principles}

From the literature and problem statement, we synthesise three core design
principles for NutriDermAI that distinguish it from prior work: 
(i)~\emph{Multimodal ABCDE computation}, where ABCD scores are computed
via fusion of image features and clinical text (not image-only);
(ii)~\emph{Structured outputs}, routing all reasoning through JSON to enable
transparency, auditability, and downstream integration (unlike black-box
models); and (iii)~\emph{Calibrated confidence}, ensuring every stage reports
uncertainty for appropriate escalation. These principles address the gap in
existing systems: most prior ABCDE-aware works compute ABCD from images alone
and ignore Evolution (E) entirely, while generic multimodal systems lack ABCDE
grounding. NutriDermAI uniquely combines both requirements in a unified
pipeline.

% ===========================================================
\section{System Overview}
% ===========================================================

The overall architecture of NutriDermAI is illustrated in
Fig.~\ref{fig:architecture}. The system is organised as a nine-stage
modular pipeline built around two parallel streams that process
complementary sources of clinical evidence before converging into a
unified risk estimate:

\begin{itemize}[leftmargin=*]
    \item \textbf{Image pipeline.} Dermatoscopic images pass through
          segmentation (SLS), ROI extraction (SLRC), and hierarchical
          classification (HC) to extract geometric, chromatic, and
          classification features encoded as floating-point vectors.
    \item \textbf{Clinical text pipeline.} Medical reports and patient
          history are processed by MedProc using clinical NLP to derive
          the Evolution~(E) score, symptom entities, and structured
          clinical features.
\end{itemize}

\textbf{Data flow and transformations.} The image pipeline produces:
(1)~a binary lesion mask from SLS (Stage~2); (2)~a tightly cropped ROI
from SLRC (Stage~3); (3)~feature embeddings, classification logits, and
raw ABCD descriptor signals (asymmetry, border, colour, diameter) from HC
(Stage~4). The text pipeline produces: a structured clinical JSON
containing extracted entities, temporal mentions, symptom-to-evolution
mappings, and an Evolution score from MedProc (Stage~5). Both pipelines
deliver their outputs—feature vectors and structured JSON—to the IT-Fusion
module (Stage~6).

\textbf{Multimodal fusion and structured JSON generation.} Stage~6
receives complementary signals from both pipelines. Image-derived features
alone cannot capture longitudinal changes or family history. Conversely,
text alone cannot assess lesion geometry, borders, or colour patterns. By
fusing both modalities, the system produces a richer evidence base. Within
Stage~6, raw descriptor signals are normalised, calibrated ABCD scores
(0--1 each) are computed via learned mappings, and these are combined with
the Evolution score from clinical text to produce a unified \emph{Clinical
Knowledge JSON}. This JSON contains: (i)~calibrated ABCDE values,
(ii)~malignancy risk probability and Total Dermatoscopic Score (TDS),
(iii)~detailed evidence summaries documenting which image and text features
contributed to each score, and (iv)~confidence levels and flags for
low-certainty predictions requiring escalation.

\textbf{VQA and conversational integration.} The Clinical Knowledge JSON
from Stage~6 becomes the authoritative structured representation consumed
by downstream modules. Stage~7 (Visual Question Answering) receives:
(1)~the cropped ROI image (224$\times$224 pixels); (2)~the user's natural-language
question ($\leq$128 tokens); (3)~the Clinical Knowledge JSON containing
ABCDE scores, risk assessment, and evidence. A vision--language model
integrates these three inputs to ground answers in ABCDE-structured evidence,
minimising hallucination and ensuring clinical alignment. The VQA output—
together with conversation history and retrieved medical guidelines—feeds
into Stage~8 (Conversational LLM with RAG), which synthesises a multi-turn
dialogue.

\textbf{Logging, audit, and deployment infrastructure.} The Clinical
Knowledge JSON and intermediate stage outputs (masks, ROI crops, confidence
scores) are logged to persistent storage for auditability and downstream
analysis. This structured, traceable pipeline is designed for secure
containerised deployment via REST API (Stage~9), with role-based access
control and encryption for sensitive medical data.

\textbf{Design rationale.} The modular, stage-wise organisation
reflects three key design principles. First, \emph{independence}: each
stage exposes well-defined inputs and outputs, allowing independent
evaluation on standard benchmarks and replacement of individual
components without disrupting the remainder of the pipeline. Second,
\emph{calibration}: every stage produces confidence scores alongside
its primary output, enabling the IT-Fusion module to weight evidence
appropriately and to flag low-confidence predictions for clinical
escalation. Third, \emph{interpretability}: by routing all downstream
reasoning through a structured Clinical JSON rather than raw model
activations, NutriDermAI provides traceable, human-readable
explanations that align with the ABCDE framework familiar to
clinicians.

\begin{figure*}[tp]
\centering
\includegraphics[width=\textwidth, height=0.82\textheight, keepaspectratio]{archi.pdf}
\caption{NutriDermAI system architecture. Solid boxes denote implemented
or planned stages; dashed boxes denote future stages requiring detailed
design. \textbf{Image pipeline} (SLS $\to$ SLRC $\to$ HC): Extracts
geometric and chromatic features. \textbf{Text pipeline} (MedProc): 
Derives Evolution~(E) score and clinical entities. \textbf{Fusion} 
(IT-Fusion): Computes calibrated ABCDE scores and produces Clinical 
Knowledge JSON. \textbf{VQA}: Grounds answers in structured evidence. 
\textbf{Conversational layer} (Stage~8): Integrates four components---
(1)~Retriever fetches relevant clinical guidelines from FAISS/KB,
(2)~Chat Memory maintains conversation history, (3)~Prompt Builder 
combines query, JSON, VQA output, and retrieved documents, and (4)~LLaMA 
generates evidence-grounded responses. \textbf{Deployment} (Stage~9): 
React frontend, FastAPI backend, and secure storage.}
\label{fig:architecture}
\end{figure*}

% ===========================================================
\section{Pipeline Architecture and Stage Design}
% ===========================================================

This section describes the design and implementation of each pipeline
stage. Stages~2--5 are fully implemented and validated; Stages~1 and
6--7 are in progress; Stages~8--9 are planned future work.

\subsection{End-to-End Workflow and Data Flow}

The overall system operates as follows:

\begin{enumerate}[leftmargin=*]
    \item \emph{Patient provides input}: Dermatoscopic image, medical history or
          symptom text, and a clinical query.
    \item \emph{Preprocessing} (Stage~1): Standardises image format, normalises
          intensities, and tokenises text.
    \item \emph{Image pipeline} (Stages~2--4): Segmentation isolates the lesion;
          ROI cropping extracts the region; classification extracts visual feature
          embeddings and raw ABCD descriptors (asymmetry, border, colour, diameter).
    \item \emph{Text pipeline} (Stage~5): Clinical NLP extracts symptoms, temporal
          mentions, and the Evolution (E) score.
    \item \emph{Multimodal fusion} (Stage~6): Combines image and text features to
          compute calibrated ABCDE values, risk scores, and a Clinical Knowledge JSON.
    \item \emph{Visual reasoning} (Stage~7): VQA module grounds the answer in both
          the ROI image and Clinical Knowledge JSON.
    \item \emph{Conversational reasoning} (Stage~8): Prompt Builder combines VQA
          output, Chat Memory history, and retrieved medical guidelines. The
          Conversational LLM integrates all signals to generate a clinically
          grounded, context-aware response.
    \item \emph{Final response}: System outputs answer, clinical explanation, and
          supporting evidence through the chatbot interface.
\end{enumerate}

This modular, stage-wise design ensures that each component can be independently
validated, replaced, or improved without disrupting the overall pipeline. All
intermediate outputs are structured (JSON or feature vectors) with explicit
confidence scores, enabling transparency and audit trails suitable for clinical
governance.

\subsection{Stage 1: Input Preprocessing}

Stage~1 performs format validation, metadata extraction, and
centralised normalisation of all inputs before branching into the image
and text pipelines. Implementation is in progress.

\subsection{Stage 2: Skin Lesion Segmentation (SLS)}

Given a full dermatoscopic image, Stage~2 produces a pixel-wise lesion
mask using ARC-UNet. Preprocessing resizes images to
512$\times$512 with bilinear interpolation, normalises intensities
using ImageNet statistics, and optionally removes hair artefacts via
morphological operations.

\subsection{Stage 3: Lesion ROI Cropping (SLRC)}

Stage~3 applies the SLS mask to extract a tightly cropped ROI via
connected-component labelling and axis-aligned bounding box computation
with 10\% padding, resized to 224$\times$224.

\subsection{Stage 4: Hierarchical Classification and Feature Extraction (HC)}

Stage~4 performs dual-headed binary and multi-class classification to
categorise lesions as benign/malignant and into one of seven dermatological
categories. An EfficientFormerV2-S2 backbone processes the ROI and
extracts high-dimensional visual feature embeddings. In parallel, four
feature extraction sub-pipelines compute intermediate geometric and
chromatic descriptors: (A)~asymmetry coefficients via mask symmetry
analysis, (B)~border irregularity metrics from contour complexity measures,
(C)~colour distribution statistics via LAB-space clustering, and (D)~lesion
diameter estimates from pixel bounding boxes. These raw geometric and
chromatic signals are passed forward as feature vectors. The actual
computation of calibrated ABCD scores and risk-weighted combinations
occurs in Stage~6 (IT-Fusion), where image-derived features are aligned
with clinical text features and normalised.

\subsection{Stage 5: Medical Information Processing (MedProc)}

Stage~5 processes patient-provided medical history and uploaded reports
(PDF, $\leq$5~MB) to derive the Evolution~(E) score and structured
clinical features. The processing pipeline applies:
\begin{itemize}[leftmargin=*]
    \item OCR for PDF reports followed by boilerplate removal and spell
          correction;
    \item tokenisation with a 512-token limit and stride-100 chunking
          for long documents; and
    \item Bio\_ClinicalBERT~\cite{alsentzer2019_bioclinicalbert} fine-tuned for named entity
          recognition (NER) and relation extraction, producing a JSON output
          containing symptoms, recommendations, and the Evolution~(E) value.
          The fine-tuned checkpoint is saved at
          \texttt{MedProc/checkpoints/medproc\_best.pt} and shared with Stage~6.
\end{itemize}

\subsection{Stage 6: Multimodal IT-Fusion and Clinical JSON Generation}

Stage~6 is the critical fusion layer that synthesises HC and MedProc
outputs into a unified, clinically interpretable representation.

\textbf{Dataset preparation.} The IT-Fusion dataset is prepared by
\texttt{dataset.ipynb}, which processes three complementary sources:
(i)~\textbf{Derm1M} (image--caption pairs aligned to a dermatology
ontology); (ii)~\textbf{MM-SkinQA} (multimodal skin QA with clinical
captions); and (iii)~\textbf{PAD-UFES-20} (tabular clinical metadata
with biopsy-confirmed labels). For each sample, two pre-trained feature
extractors are applied: (a)~the \textbf{Bio\_ClinicalBERT} encoder
(\texttt{emilyalsentzer/Bio\_ClinicalBERT}), fine-tuned in Stage~5
and loaded from the MedProc checkpoint, extracts 768-dim text
embeddings; and (b)~the \textbf{EfficientFormerV2-S2 backbone},
fine-tuned in Stage~4 (HC) and exported via \texttt{hc\_model.py}
(\texttt{load\_hc\_backbone()}), extracts 288-dim image features.
The two encoders are explicitly reused from their respective training
stages---not retrained---ensuring consistent representation alignment
across the pipeline.

\textbf{Fusion architecture.} The stage receives: (i)~feature vectors
and raw ABCD descriptor signals from Stage~4 (HC); (ii)~classification
logits (benign/malignant and 8-class); (iii)~the Evolution~(E) score
and clinical symptom features from Stage~5 (MedProc); and
(iv)~confidence estimates for each signal. The fusion process
comprises four key operations: \emph{(1)~Calibration}: numerical
features are z-normalised using training-set statistics, and confidence
scores are recalibrated via Platt scaling. \emph{(2)~ABCD Score
Computation}: raw descriptor signals from Stage~4 are mapped to
clinically calibrated ABCD scores (0--1 each) using learned mappings
trained on clinician-annotated data. \emph{(3)~Risk Scoring}: ABCDE
scores are combined via a weighted ensemble to produce a final
malignancy risk probability and a Total Dermatoscopic Score (TDS).
Conflict resolution (e.g., when image suggests benign but text
suggests evolution) is handled via confidence-weighted averaging and
flagged for clinician review. \emph{(4)~JSON Serialisation}: a
\emph{Clinical Knowledge JSON} is generated containing ABCDE values,
risk score, evidence summaries, and a ranking of which signals (image
vs.\ text) most contributed to the final decision. This JSON serves as
the authoritative structured representation used by downstream VQA,
RAG, and audit systems. This stage is currently in progress.

\subsection{Stage 7: Visual Question Answering (VQA)}

Stage~7 converts the fused Clinical Knowledge JSON into natural-language
answers grounded in structured ABCDE evidence. It receives three inputs:
(1)~the cropped ROI image (224$\times$224 pixels) from Stage~3; (2)~the
user's natural-language query ($\leq$128 tokens); and (3)~the Clinical
Knowledge JSON from Stage~6 containing calibrated ABCDE scores, risk
assessment, and evidence rankings. An R-LLaVA-style vision--language
model jointly attends to image, text, and structured JSON to produce:
a natural-language answer addressing the user's question, an evidence
summary documenting which ABCDE signals and clinical facts support the
answer, and a confidence estimate. By grounding responses in the
structured JSON rather than raw image features, VQA minimises hallucination
and ensures clinical alignment. VQA is invoked only after IT-Fusion produces
a consolidated JSON, preserving explainability and audit trails. This stage
is currently in progress.

\subsection{Stage 8: Conversational LLM with RAG (Planned)}

To support multi-turn clinical dialogue, Stage~8 integrates four key
components: a \emph{Retriever}, a \emph{Prompt Builder}, a \emph{Chat Memory},
and a \emph{Conversational LLM}. Operation is as follows:

\begin{enumerate}[leftmargin=*]
    \item \emph{Retriever (FAISS / Knowledge Base)}: Given the current user
          query, a FAISS vector database retrieves the top-$k$ most relevant
          documents from a curated knowledge base containing ABCDE clinical
          rules, dermatology guidelines, treatment protocols, and evidence-based
          recommendations. Retrieval operates over both the query text and the
          Clinical Knowledge JSON to ensure context-aware document selection.
    \item \emph{Chat Memory}: Maintains a persistent conversation history
          throughout a multi-turn dialogue session. Each query--response pair is
          logged with timestamps and clinical context (lesion ID, patient
          identifier if applicable) for auditability and longitudinal tracking.
    \item \emph{Prompt Builder}: Constructs a unified, clinically-aligned
          prompt that integrates: (a)~the current user query, (b)~the Clinical
          Knowledge JSON from Stage~6 (ABCDE scores, risk assessment, evidence),
          (c)~VQA output from Stage~7 (answer and evidence summary), (d)~retrieved
          documents from the Retriever, and (e)~relevant conversation history from
          Chat Memory. The prompt is formatted to encourage reasoning transparency
          and evidence-grounded responses.
    \item \emph{Conversational LLM (LLaMA-based)}: Receives the constructed
          prompt and generates a natural-language response that addresses the
          user's question while maintaining clinical alignment. The LLM output
          includes the answer, supporting evidence, and any recommended follow-up
          actions or escalation flags.
\end{enumerate}

This stage is planned as future work and will undergo rigorous stress testing
for hallucination, out-of-scope questions, and adversarial prompts before
clinical deployment.

\subsection{Stage 9: Frontend, Backend, and Deployment (Planned)}

Stage~9 will expose the system as a chat-based web application
comprising: a React frontend for image uploads, report attachment, and
chat; a FastAPI backend orchestrating Stages~2--8 as containerised
microservices; and secure artifact storage (masks, JSON, logs) for
auditability. Deployment will include continuous monitoring for drift,
error patterns, and user feedback, informed by dermatology chatbot
design
guidelines~\cite{dermatologychatbot2024,casestudy2024_wizard_of_oz,
amie2025,convagents2025_effectiveness}.

% ===========================================================
\section{Stage Overview}
% ===========================================================
\label{sec:stage_overview}

Table~\ref{tab:stage_overview} summarises all pipeline stages,
consolidating inputs, datasets, preprocessing, models, implementation
status, and key metrics. Stages~2--5 are fully implemented and
validated on standard benchmarks; Stages~1 and 6--7 are in progress.

\begin{table*}[tp]
\centering
\footnotesize
\renewcommand{\arraystretch}{1.35}
\caption{NutriDermAI Pipeline Stage Overview}
\label{tab:stage_overview}
\setlength{\tabcolsep}{4pt}
\begin{tabular}{
  c
  L{1.9cm}
  L{2.3cm}
  L{3.0cm}
  L{2.2cm}
  C{1.1cm}
  L{2.2cm}
}
\toprule
\textbf{Stage} &
\textbf{Input / Output} &
\textbf{Dataset (Train/Val/Test)} &
\textbf{Preprocessing} &
\textbf{Model} &
\textbf{Status} &
\textbf{Key Metric} \\
\midrule

1 &
Raw image / text &
N/A &
Format validation, metadata extraction &
Centralised normaliser &
In Prog. &
--- \\

\midrule

\multirow{2}{*}{2 (SLS)} &
\multirow{2}{*}{Image $\to$ Mask} &
\multirow{2}{*}{\makecell[l]{ISIC 2016--17\\(2.9K/400/979)}} &
\multirow{2}{*}{\makecell[l]{Resize 512$\times$512,\\norm., hair removal}} &
\multirow{2}{*}{ARC-UNet} &
\multirow{2}{*}{\checkmark} &
Dice: 0.8506 \\
& & & & & & Bdry F1: 0.8815 \\

\midrule

3 (SLRC) &
Mask $+$ Image $\to$ ROI &
ISIC 2016--17 (2.9K/400/979) &
Morphology, connected components, bbox, 10\% pad &
OpenCV 4.5 &
\checkmark &
Coverage: 98.4\% \\

\midrule

\multirow{2}{*}{4 (HC)} &
\multirow{2}{*}{\makecell[l]{ROI $\to$\\Label $+$ Features}} &
\multirow{2}{*}{\makecell[l]{ISIC $+$ HAM $+$\\ISIC 19--24\\(35K/5K/10K)}} &
\multirow{2}{*}{\makecell[l]{Augmentation,\\normalisation, MixUp}} &
\multirow{2}{*}{\makecell[l]{EfficientFormer-\\V2-S2}} &
\multirow{2}{*}{\checkmark} &
Acc: 94.7\% (bin.) \\
& & & & & & Acc: 89.2\% (7-cls.) \\

\midrule

5 (MedProc) &
Report $\to$ Symptoms $+$ E &
MIMIC-III/IV (50K/5K/10K) &
OCR, cleanup, tokenisation ($\leq$512 tokens, stride 100) &
Bio\_ClinicalBERT~\cite{alsentzer2019_bioclinicalbert} &
\checkmark &
NER F1: 0.87--0.89 \\

\midrule

6 (Fusion) &
HC $+$ MedProc $\to$ Clinical JSON &
\makecell[l]{Derm1M\\MM-SkinQA\\PAD-UFES-20} &
Feature alignment, z-norm, embedding &
MARIA Transformer &
In Prog. &
Risk AUC (TBD) \\

\midrule

7 (VQA) &
JSON $+$ ROI $+$ Q $\to$ Answer &
DermaVQA $+$ PathVQA (80K/10K/15K) &
Prompt conversion, tokenisation &
R-LLaVA &
In Prog. &
VQA F1 (TBD) \\

\bottomrule
\end{tabular}
\end{table*}

% ===========================================================
\section{Experimental Results}
% ===========================================================

\subsection{Experimental Setup and Dataset Justification}

\textbf{Datasets and selection rationale.} We employ multiple public
datasets, each chosen for specific strengths: (i)~\textbf{ISIC
2016--2017} (4,279 images): Gold-standard dermatoscopy dataset with
expert-annotated lesion masks, used for segmentation (SLS stage);
(ii)~\textbf{HAM10000} (10,015 images): Large-scale, ethnically diverse
lesion classification dataset covering 7 categories and multiple skin
tones, used for classification (HC stage); (iii)~\textbf{ISIC
2018--2024} (50,000+): Extended classification benchmark with updated
annotations, used for HC stage fine-tuning; (iv)~\textbf{MIMIC-III/IV}:
Large-scale medical NLP corpus (60K+ documents) with longitudinal patient
records, used for MedProc stage training and Evolution~(E) score
calibration; (v)~\textbf{Derm1M}: Emerging large-scale multimodal
dermatology dataset (1M+ image--text pairs) for multimodal fusion research.

\textbf{Hardware, software, and implementation.} All experiments are
conducted on an NVIDIA A100 GPU (40~GB memory) with 48 CPU cores and
512~GB RAM. Image stages (SLS, SLRC, HC) are implemented in PyTorch
2.1 with OpenCV for image processing. Clinical NLP (MedProc) uses HuggingFace Transformers with
Bio\_ClinicalBERT~\cite{alsentzer2019_bioclinicalbert} fine-tuned for
clinical NER and Evolution scoring. All code is containerised using
Docker for reproducibility.

\textbf{Hyperparameter configurations and training details.} Key training
configurations are detailed below:

\begin{itemize}[leftmargin=*]
    \item \textbf{SLS (ARC-UNet)}: Batch size 8, learning rate $10^{-4}$,
          optimizer Adam, loss hybrid Dice+CE (ratio 0.6/0.4), 100 epochs,
          early stopping patience 10, $L_2$ regularisation $\lambda=10^{-5}$.
          Augmentation: random flips, rotations $\pm 15°$, elastic deformations.
    \item \textbf{HC (EfficientFormerV2-S2)}: Batch size $16 \times 2$
          (gradient accumulation), learning rate $10^{-3}$ (head)~/$10^{-4}$
          (backbone), OneCycleLR scheduler, loss AsymmetricFocalLoss
          ($\gamma=2.0$), 80 epochs, weighted sampling for class balance,
          MixUp/CutMix augmentation.
    \item \textbf{MedProc (Bio\_ClinicalBERT)}: Batch size 32, learning rate $5 \times 10^{-5}$,
          optimizer AdamW, 10 epochs, 512-token limit with stride-100 chunking,
          dropout 0.1, label smoothing 0.1. Evaluation uses ECE (Expected
          Calibration Error) and per-entity F1 scores.
\end{itemize}

We present stage-wise results for the four completed stages (2--5).
For each stage, we document the dataset, algorithm, training
configuration, and measured performance metrics.

\subsection{Stage 2: Skin Lesion Segmentation --- ARC-UNet}

\textbf{Dataset.} ISIC 2016--2017: 4,279 images (2,900 train / 400 val
/ 979 test).

\textbf{Algorithm.} ARC-UNet with squeeze-and-excitation attention,
residual connections, and an asymmetric encoder--decoder.

\textbf{Training.} Hybrid Dice$+$CE loss (ratio 0.6/0.4), Adam
($\mathrm{LR}=10^{-4}$), batch size 8, 100 epochs, L$_2$
regularisation ($\lambda=10^{-5}$), early stopping (patience 10),
augmentation (flips, rotations $\pm15^{\circ}$, elastic deformations).

\textbf{Results.} Table~\ref{tab:sls_results} reports single-pass and
7-view test-time augmentation (TTA) performance. ARC-UNet achieves
Boundary F1 of 0.8815 with balanced sensitivity and specificity.
Single-pass inference is recommended for clinical deployment.

\begin{table}[t]
\centering
\renewcommand{\arraystretch}{1.2}
\caption{Stage 2 (SLS) --- ARC-UNet Segmentation Results}
\label{tab:sls_results}
\setlength{\tabcolsep}{4pt}
\begin{tabular}{l S[table-format=1.4] S[table-format=1.4] r}
\toprule
\textbf{Metric} & {\textbf{Single-Pass}} & {\textbf{TTA (7-view)}} & \textbf{Gain} \\
\midrule
Dice Coefficient        & 0.8506 & 0.8512 & $+$0.07\% \\
IoU                     & 0.7718 & 0.7726 & $+$0.10\% \\
Pixel Accuracy (\%)     & 93.70  & 93.75  & $+$0.05\% \\
Boundary F1-Score       & 0.8801 & 0.8815 & $+$0.16\% \\
Sensitivity (Recall)    & 0.9102 & 0.9118 & $+$0.18\% \\
Specificity             & 0.9756 & 0.9761 & $+$0.05\% \\
\midrule
\multicolumn{4}{l}{\textit{Optimal threshold: 0.44 (well-calibrated)}} \\
\bottomrule
\end{tabular}
\end{table}

\subsection{Stage 3: Lesion ROI Cropping --- OpenCV Pipeline}

\textbf{Dataset.} Derived from ISIC 2016/2017 SLS output masks
(4,279 images; no separate dataset).

\textbf{Algorithm.} Deterministic four-step pipeline:
\begin{enumerate}[leftmargin=*,topsep=2pt,itemsep=1pt]
    \item \emph{Morphological cleanup:} binary closing (disk, $r=3$) to
          fill segmentation holes.
    \item \emph{Connected-component labelling:} 8-connectivity; retain
          the largest component.
    \item \emph{Bounding box extraction:} axis-aligned bounding box with
          10\% padding on all sides.
    \item \emph{Aspect-ratio-preserving resize:} bilinear interpolation
          to 224$\times$224 with black-border padding if needed.
\end{enumerate}

\textbf{Results.} Table~\ref{tab:slrc_results} confirms high lesion
coverage (98.4\%) and border preservation (96.7\%) at 12.3~ms latency
per ROI.

\begin{table}[t]
\centering
\renewcommand{\arraystretch}{1.2}
\caption{Stage 3 (SLRC) --- ROI Cropping Results}
\label{tab:slrc_results}
\setlength{\tabcolsep}{5pt}
\begin{tabular}{L{3.8cm} L{3.4cm}}
\toprule
\textbf{Metric} & \textbf{Value} \\
\midrule
Mean Lesion Coverage (\% pixels retained)    & 98.4\% $\pm$ 1.2\% \\[2pt]
Border Preservation (\% contour within ROI)  & 96.7\% $\pm$ 2.1\% \\[2pt]
Mean ROI Size Variation (vs.\ anat.\ diam.)  & 0.34~mm \newline (IQR: 0.18--0.65~mm) \\[2pt]
Computational Time per ROI (CPU)             & 12.3~ms \\
\bottomrule
\end{tabular}
\end{table}

\subsection{Stage 4: Hierarchical Classification --- EfficientFormerV2-S2}

\textbf{Dataset.} ISIC 2018 $+$ HAM10000 $+$ ISIC 2019--2024; 50,000
balanced samples (70/10/20 split, no patient overlap).

\textbf{Algorithm.} EfficientFormerV2-S2 with dual-head classification
(benign/malignant $+$ 7-class) and parallel ABCD sub-pipelines.

\textbf{Training.} AsymmetricFocalLoss ($\gamma=2.0$), OneCycleLR
(head $10^{-3}$, backbone $10^{-4}$), batch
$16{\times}2$ gradient accumulation, 80 epochs, MixUp/CutMix,
WeightedRandomSampler.

\textbf{Results.} Table~\ref{tab:hc_results} shows that explicit ABCD
features improve both interpretability and performance. TTA provides
approximately 1\% improvement for critical cases; single-pass is
acceptable for routine triage.

\begin{table}[h]
\centering
\small
\renewcommand{\arraystretch}{1.2}
\caption{Stage 4 (HC) --- Hierarchical Classification Results}
\label{tab:hc_results}
\setlength{\tabcolsep}{5pt}
\begin{tabular}{l c c c}
\toprule
\textbf{Metric} & \textbf{Single-Pass} & \textbf{TTA (7-view)} & \textbf{Gain} \\
\midrule
\multicolumn{4}{l}{\textit{Binary Classification (Benign / Malignant)}} \\
\midrule
Accuracy (\%)        & 94.7   & 95.3   & $+$0.6\%  \\
Sensitivity          & 0.9512 & 0.9618 & $+$1.06\% \\
Specificity          & 0.9387 & 0.9425 & $+$0.38\% \\
AUC-ROC              & 0.9751 & 0.9803 & $+$0.52\% \\
\midrule
\multicolumn{4}{l}{\textit{Multi-Class Classification (7 Categories)}} \\
\midrule
Accuracy (\%)              & 89.2   & 90.1   & $+$0.9\%  \\
Weighted F1-Score          & 0.8876 & 0.8979 & $+$1.03\% \\
Macro F1-Score             & 0.8654 & 0.8762 & $+$1.08\% \\
Per-Class Balanced Acc.    & 0.8905 & 0.8998 & $+$0.93\% \\
\bottomrule
\end{tabular}
\end{table}

\subsection{Stage 5: Medical Information Processing --- BioBERT}

\textbf{Dataset.} MIMIC-III/IV: 50K train / 5K val / 10K test (no
patient overlap).

\textbf{Algorithm.} Bio\_ClinicalBERT~\cite{alsentzer2019_bioclinicalbert}
augmented with two transformer blocks for clinical NER, relation
extraction, and rule-based Evolution~(E) scoring.

\textbf{Training.} AdamW ($\mathrm{LR}=5{\times}10^{-5}$), batch 32,
10 epochs, 512-token limit with stride-100 chunking, dropout 0.1,
label smoothing 0.1.

\textbf{Results.} Table~\ref{tab:medproc_results} shows NER F1 of
0.87--0.89 across clinical entity types, with well-calibrated
uncertainty (ECE 0.0201), suitable for downstream fusion.

\begin{table}[h]
\centering
\small
\renewcommand{\arraystretch}{1.2}
\caption{Stage 5 (MedProc) --- Clinical NLP Results}
\label{tab:medproc_results}
\setlength{\tabcolsep}{4pt}
\begin{tabular}{l c c c}
\toprule
\textbf{Task / Entity} & \textbf{Prec.} & \textbf{Rec.} & \textbf{F1} \\
\midrule
\multicolumn{4}{l}{\textit{Named Entity Recognition (NER)}} \\
\midrule
Symptoms            & 0.8901 & 0.8596 & 0.8742 \\
Body sites          & 0.9087 & 0.8796 & 0.8934 \\
Temporal mentions   & 0.8312 & 0.8011 & 0.8156 \\
Severity indicators & 0.7956 & 0.7698 & 0.7821 \\
\midrule
\multicolumn{4}{l}{\textit{Relation Extraction (Symptom-to-Evolution)}} \\
\midrule
Progression detected (P / R) & 0.8764 & 0.8523 & --- \\
AUC-ROC              & \multicolumn{3}{c}{0.8934} \\
\midrule
\multicolumn{4}{l}{\textit{Evolution (E) Score Quality}} \\
\midrule
Expected Calibration Error (ECE) & \multicolumn{3}{c}{0.0201} \\
Mean Absolute Error (scale 0--1) & \multicolumn{3}{c}{0.1547} \\
\bottomrule
\end{tabular}
\end{table}

\subsection{Integration and Data Flow Validation}

End-to-end integration tests across Stages~2--5 confirm correct data
flow: SLS mask (512$\times$512 binary) $\to$ SLRC ROI
(224$\times$224 RGB) $\to$ HC classification logits and ABCD features
$+$ MedProc JSON (symptoms and E-score). Structured outputs---JSON for
MedProc, numpy arrays for image stages---are properly serialised and
deserialised between stages. Per-sample latency on an NVIDIA A100 is:
SLS 280~ms, SLRC 12~ms, HC 45~ms, MedProc 180~ms; total
\textbf{517~ms}, which is acceptable for clinical batch processing.

% ===========================================================
\section{Discussion}
% ===========================================================

\subsection{Validation of Modular Design and Component Performance}

The four-stage modular architecture yields several practical advantages.
First, each stage can be evaluated independently on established
benchmarks (ISIC for SLS/HC; MIMIC for MedProc; deterministic
validation for SLRC), facilitating transparent reporting and
reproducibility. SLS achieves Dice 0.8506 and Boundary F1 0.8815,
competitive with recent state-of-the-art U-Net variants. HC attains
94.7\% binary accuracy and 89.2\% 7-class accuracy with strong
per-class balance, demonstrating that explicit ABCD feature extraction
does not compromise classification performance. MedProc achieves NER F1
of 0.87--0.89 with low calibration error (ECE 0.0201), suitable for
downstream fusion. Together, these results confirm that modular
independence does not compromise individual component quality.

Second, SLS and HC are deployable as standalone components, with MedProc
available as an optional enrichment layer. This flexibility supports
iterative deployment: institutions may begin with image-only analysis
(SLS~+~HC) and incrementally integrate text processing as clinical
workflows mature. Third, all outputs are structured (JSON or numpy
arrays) with calibrated confidence scores, enabling explainability and
audit trails. Fourth, the microservice-compatible design supports
containerised cloud and on-premise deployment with role-based access
control.

\subsection{Implications of ABCDE Computation in IT-Fusion}

A key architectural choice is to compute ABCD scores in Stage~6
(IT-Fusion) rather than Stage~4 (HC). This design decision reflects two
insights: (i)~ABCD metrics cannot be properly calibrated using image
features alone, as clinician-annotated ABCDE labels typically reflect
integrated judgement combining visual, temporal, and contextual factors;
and (ii)~Evolution (E) depends exclusively on clinical text and cannot
be inferred from images. By deferring ABCD score computation to the
fusion stage, NutriDermAI ensures that all five criteria are computed
consistently and that the final risk estimate properly integrates both
modalities. Future work may explore learning-based ABCD calibration
mappings via annotation of a prospective cohort.

\subsection{Safety and Ethical Considerations}

NutriDermAI is designed as a decision support tool, not a standalone
diagnostic system. All outputs include explicit confidence levels, and
low-confidence predictions automatically trigger escalation pathways to
human dermatologists. The structured Clinical JSON audit trail documents
which image and text features contributed to each decision, enabling
post-hoc review and quality assurance.

Deployment must comply with emerging regulatory frameworks for AI-based
medical devices. The FDA's recent guidance on AI/ML-based Software as a
Medical Device (SaMD)~\cite{fda2021_aiml_samd} emphasises transparency,
performance monitoring, and post-market surveillance. The EU AI Act
classifies medical AI as high-risk, requiring documented risk assessments,
data governance, and human oversight~\cite{euaiact2024}. NutriDermAI's
modular design and structured outputs facilitate compliance: each stage
can be independently validated, and the JSON evidence trail enables
detailed audits.

Fairness and bias are critical concerns. The current training cohorts
(ISIC, HAM10000) are known to skew towards lighter skin tones, potentially
reducing performance on darker skin phenotypes~\cite{kinyanjui2021_color_bias_skin,
buolamwini2019_gender_shades}. Prospective fairness audits across skin
tones, anatomical sites, and lesion types are planned. If disparities are
detected, targeted data augmentation, fine-tuning on diverse cohorts, and
ensemble methods can mitigate bias. Multinational clinical validation with
diverse patient populations is essential before broad deployment.

\subsection{Key Strengths and Architectural Contributions}

The NutriDermAI architecture demonstrates several distinct strengths that
differentiate it from existing dermatology AI systems:

\begin{itemize}[leftmargin=*]
    \item \textbf{Multimodal integration.} By combining dermatoscopic images,
          clinical text, and user queries, the system captures information that
          no single modality can provide alone. This fusion enables more robust
          and contextually grounded clinical reasoning.
    \item \textbf{Structured explainability.} The Clinical Knowledge JSON serves
          as an intermediate representation that makes all reasoning transparent.
          Clinicians can audit which image and text features contributed to each
          ABCDE score and risk decision, supporting trust and compliance.
    \item \textbf{ABCDE-aware reasoning.} Unlike generic medical VQA systems, our
          approach explicitly computes and reasons over the ABCDE rule---the
          clinical gold standard for melanoma risk assessment. This alignment with
          dermatology practice standards improves clinical relevance and uptake.
    \item \textbf{Knowledge-grounded conversational AI.} By integrating Retrieval-Augmented
          Generation (RAG), the system grounds responses in verified medical guidelines
          and dermatology evidence, reducing hallucinations and improving accuracy
          compared to LLM-only approaches.
    \item \textbf{Modular design for iterative deployment.} Each stage is independently
          validated and can be improved or replaced independently. Institutions can
          start with image-only analysis (SLS~+~HC) and progressively integrate text
          processing and conversational reasoning as clinical workflows mature.
    \item \textbf{Calibrated confidence and clinical escalation.} Every stage reports
          uncertainty estimates; low-confidence predictions automatically trigger
          escalation pathways to human dermatologists, ensuring patient safety.
\end{itemize}

Four core modules are fully implemented and evaluated:
\begin{itemize}[leftmargin=*]
    \item \textbf{Stage 2 (SLS):} Dice 0.8506, IoU 0.7718, 93.70\%
          pixel accuracy (Table~\ref{tab:sls_results}).
    \item \textbf{Stage 3 (SLRC):} 98.4\% coverage, 96.7\% border
          preservation, 12.3~ms latency (Table~\ref{tab:slrc_results}).
    \item \textbf{Stage 4 (HC):} 94.7\% binary accuracy (TTA 95.3\%),
          89.2\% 7-class accuracy (TTA 90.1\%)
          (Table~\ref{tab:hc_results}).
    \item \textbf{Stage 5 (MedProc):} NER F1 0.8742--0.8934, ECE
          0.0201 (Table~\ref{tab:medproc_results}).
\end{itemize}

Remaining stages requiring implementation are: Stage~1 (input
normalisation), Stage~6 (IT-Fusion multimodal transformer with ABCDE
calibration), and Stage~7 (R-LLaVA VQA). Stages~8--9 (conversational LLM
with RAG, and deployment) are future work pending completion of Stages~1
and 6--7. Integration and validation of the complete pipeline is
targeted for Q4 2026.

\subsection{Generalisability and Deployment Considerations}

The pipeline's modular design supports generalisation across settings.
ISIC benchmarks span multiple imaging devices and geographies, suggesting
reasonable device robustness. However, real-world deployment in clinical
settings will likely encounter out-of-distribution inputs (poor lighting,
large artefacts, non-standard cameras). Robust out-of-distribution
detection and failure mode analysis are planned. Furthermore, sequential
processing (SLS $\to$ SLRC $\to$ HC $\to$ MedProc $\to$ Fusion)
introduces latency. Optimisations include batch processing, model
quantisation, and edge deployment of lightweight models for resource-limited
settings.

The English-only text processing is a significant limitation for
international deployment. Multilingual clinical NLP models and translated
prompts will extend applicability. Stateless inference (processing each
case independently) also limits clinical utility; a stateful version
maintaining longitudinal ABCDE trends within a patient's record is
planned for integration with clinical EHR systems.

% ===========================================================
\section{Limitations and Future Work}
% ===========================================================

\subsection{Current Limitations}

The NutriDermAI prototype, though comprehensive in design, carries several
known limitations that inform the research roadmap:

\begin{itemize}[leftmargin=*]
    \item \textbf{English-only text processing.} All clinical NLP components
          are trained on English-language corpora. In global health contexts
          where patients may provide reports in other languages, a multilingual
          extension is necessary~\cite{pfohl2022_multilingual_clinical_nlp}.
    \item \textbf{Skin tone bias.} Like many dermatology datasets, the training
          cohorts (ISIC, HAM10000) skew towards lighter skin tones. Systematic
          fairness audits across diverse skin tones are planned~\cite{buolamwini2019_gender_shades,
          kinyanjui2021_color_bias_skin}.
    \item \textbf{Dataset curation gaps.} The current HC stage uses curated
          ISIC + HAM10000 data; in-the-wild deployment may encounter artefacts
          (hair, ink, lighting) not well-represented in training sets. Robust
          out-of-distribution detection is planned.
    \item \textbf{Stateless inference.} The current pipeline processes each image
          independently without longitudinal context. Clinical deployment requires
          temporal patient records; a stateful version maintaining historical
          ABCDE trends is planned.
    \item \textbf{Staged validation.} While Stages~2--5 have been independently
          validated, end-to-end performance of the fused system (Stages~6--7) is
          pending completion and clinical evaluation.
\end{itemize}

\subsection{Future Work and Clinical Validation Roadmap}

\begin{itemize}[leftmargin=*]
    \item \textbf{Completion of core pipeline stages.} Stage~1 (input normalisation),
          Stage~6 (IT-Fusion multimodal transformer), and Stage~7 (R-LLaVA VQA) are
          in active development. Integration and validation are expected by Q4 2026.
    \item \textbf{Prospective clinical validation.} Once Stages~1--7 are complete,
          a prospective study comparing NutriDermAI's ABCDE scores and risk rankings
          against board-certified dermatologist consensus will establish diagnostic
          concordance and clinical utility~\cite{tjoa2019_clinical_validation_ai}.
    \item \textbf{Fairness and bias audits.} Systematic evaluation across skin tones,
          lesion types, and anatomical sites will be conducted using established
          fairness metrics~\cite{kinyanjui2021_color_bias_skin,buolamwini2019_gender_shades}.
    \item \textbf{Regulatory compliance and deployment.} Compliance with FDA guidance
          on AI/ML-based Software as a Medical Device (SaMD)~\cite{fda2021_aiml_samd}
          and the EU AI Act~\cite{euaiact2024} will be undertaken. Stage~9 deployment
          will include role-based access, audit logging, data encryption, and continuous
          monitoring for model drift and failure modes.
    \item \textbf{Conversational robustness and safety.} Stage~8 (conversational LLM
          with RAG) will be stress-tested for hallucination, adversarial prompts, and
          out-of-scope questions. Guardrails following AMIE and clinical chatbot best
          practices~\cite{amie2025,dermatologychatbot2024} will be implemented.
    \item \textbf{Multi-language and cross-cultural adaptation.} Multilingual clinical
          NLP models and culturally sensitive VQA prompts will extend NutriDermAI's
          applicability to low-resource regions.
    \item \textbf{Real-time inference optimisation.} Model distillation and quantisation
          will enable deployment on resource-constrained devices (smartphones, rural
          clinics with limited bandwidth).
\end{itemize}

% ===========================================================
\section{Conclusion}
% ===========================================================

This paper presents NutriDermAI, a modular, ABCDE-aware multimodal
dermatology decision support system with four validated core stages:
skin lesion segmentation (Dice 0.8506, Boundary F1 0.8815), ROI
extraction (98.4\% coverage, 12.3~ms), hierarchical classification
(94.7\% binary accuracy, 89.2\% 7-class), and clinical NLP for
evolution scoring (NER F1 0.87--0.89, ECE 0.0201). End-to-end
integration across these stages is confirmed with a total pipeline
latency of 517~ms on an NVIDIA A100.

\textbf{Core architectural contribution.} The critical contribution of
NutriDermAI is the principled fusion of image and text modalities to
compute clinically calibrated ABCDE scores, a task typically performed
ad-hoc by dermatologists with significant inter-observer variation. By
routing all downstream reasoning through a structured Clinical Knowledge
JSON, the system provides interpretability, independent stage evaluation,
and redeployable components suitable for clinical governance and audit.
Furthermore, by integrating knowledge-grounded conversational AI via
Retrieval-Augmented Generation, NutriDermAI supports multi-turn clinical
dialogue while minimising hallucination and ensuring alignment with
verified medical evidence.

\textbf{Thesis statement.} The proposed architecture bridges perception,
reasoning, and interaction by integrating multimodal learning, structured
clinical inference, and knowledge-grounded conversational AI, thereby
enabling a scalable and explainable dermatology decision support system.
This integrated design addresses a critical gap in the literature, where
existing systems either provide image-only classification without clinical
context, or offer generic medical VQA without ABCDE-aware reasoning. By
combining all three capabilities in a unified, modular pipeline, NutriDermAI
positions itself as a comprehensive platform for both clinical decision support
and patient education.

\textbf{Path forward.} NutriDermAI provides a validated research foundation for
ABCDE-aware dermatology AI systems that combine strong empirical performance
with transparency and modular reusability. Prospective clinical validation
studies, completion of Stages~1, 6, and 7, and fairness audits across
diverse populations are the immediate priorities. If such validation
confirms diagnostic concordance with board-certified dermatologists,
NutriDermAI could serve as a scalable triage tool in resource-limited
settings where specialist access is restricted~\cite{dreremoula2022_derm_shortage},
while simultaneously supporting clinical training and reducing diagnostic
variability in high-income regions.


\bibliographystyle{IEEEtran}
\bibliography{references}

\end{document}